\documentclass[12pt,a4paper]{article}
\usepackage[utf8]{inputenc}
\usepackage[indonesian]{babel}
\usepackage{amsmath}
\usepackage{amsfonts}
\usepackage{amssymb}
\usepackage{geometry}

\geometry{margin=2.5cm}

\title{Penyelesaian Volume Benda dengan Metode Shell}
\author{Ahmad Fakhrul Bawani}
\date{}

\begin{document}

\maketitle

\section*{Soal}

Diketahui:
\begin{itemize}
  \item Fungsi pembatas: $y = 3\sqrt{x^2 + 1}$
  \item Batas-batas: $x = 0$ dan $x = \sqrt{15}$
  \item Pias: Sejajar sumbu-$y$ (diintegrasikan terhadap $x$)
  \item Sumbu putar: Sumbu-$y$
\end{itemize}
Tentukan volumenya menggunakan metode shell!

\subsection*{1. Rumus Umum}
Karena pias sejajar sumbu-$y$ dan diputar mengelilingi sumbu-$y$, fungsi diintegrasikan terhadap $x$:
\begin{equation*}
  V = 2\pi \int_{a}^{b} (\text{jari-jari}) \cdot (\text{tinggi}) \, dx
\end{equation*}
Dimana jari-jari tabung (jarak pias ke sumbu putar) adalah $x$ dan tinggi tabung adalah $y = 3\sqrt{x^2 + 1}$.

\subsection*{2. Solusi Integrasi}
Substitusikan komponen yang diketahui ke dalam rumus:
\begin{align*}
  V &= 2\pi \int_{0}^{\sqrt{15}} x \cdot \left(3\sqrt{x^2 + 1}\right) \, dx \\
  V &= 6\pi \int_{0}^{\sqrt{15}} x(x^2 + 1)^{\frac{1}{2}} \, dx
\end{align*}

Gunakan metode substitusi integral:
\begin{align*}
  \text{Misalkan: } & u = x^2 + 1 \implies du = 2x \, dx \implies x \, dx = \frac{1}{2} \, du \\
  \text{Batas baru: } & \text{Untuk } x = 0 \implies u = 0^2 + 1 = 1 \\
  & \text{Untuk } x = \sqrt{15} \implies u = (\sqrt{15})^2 + 1 = 16
\end{align*}

Substitusikan variabel baru $u$ ke dalam integral:
\begin{align*}
  V &= 6\pi \int_{1}^{16} u^{\frac{1}{2}} \cdot \left(\frac{1}{2} \, du\right) \\
  V &= 3\pi \int_{1}^{16} u^{\frac{1}{2}} \, du \\
  V &= 3\pi \left[ \frac{2}{3}u^{\frac{3}{2}} \right]_{1}^{16} \\
  V &= 2\pi \left[ u\sqrt{u} \right]_{1}^{16} \\
  V &= 2\pi \left( \left[ 16\sqrt{16} \right] - \left[ 1\sqrt{1} \right] \right) \\
  V &= 2\pi (16 \cdot 4 - 1) \\
  V &= 2\pi (64 - 1) \\
  V &= 2\pi (63) \\
  V &= 126\pi
\end{align*}

Jadi, volume benda putar tersebut adalah \textbf{$126\pi$ satuan volume}.

\end{document}